% SPDX-License-Identifier: AGPL-3.0-or-later
% A symbol no font has: the segment modality, one glyph drawn in FontForge and
% saved as an OpenType font, SegmentSymbol.otf, at U+01A7 ("Ƨ", the reversed S
% it resembles, so that copying it gives a real character). The build puts the
% font next to the document, so it is loaded by its file name; declared as a
% symbol font, it takes the script sizes from the math machinery.
\DeclareFontFamily{TU}{segmentsymbol}{}
\DeclareFontShape{TU}{segmentsymbol}{m}{n}{<->"[SegmentSymbol.otf]"}{}
\DeclareSymbolFont{segmentsymbols}{TU}{segmentsymbol}{m}{n}
% An ordinary math symbol, like \exists: \segment_X takes its variable as
% a subscript.
\Umathchardef\segment="0 \symsegmentsymbols "01A7\relax
% Character protrusion and font expansion, as in the PDF: the viewer
% protrudes and expands as the document says, and only then.
\usepackage{microtype}
